\section{Introduction}

The image of the \textit{life ray} did not begin as an academic construct. It began in a letter.

Years ago, I wrote a personal letter to someone whose presence had unexpectedly marked my own life.
In that letter I tried, almost helplessly, to describe what it means for one person to cross the path of another.
I did not have formal language for it at the time. I reached instead for metaphors: parallel timelines, points of contact,
a shared space where two lives temporarily coil around each other like strands of a helix.
I spoke of ``fragments'' -- memory particles left behind in us by others.

Only later did I understand that this attempt to describe a single human encounter was already the seed of a theory.

The present text is an attempt to formalize that intuition. I propose a conceptual framework that I call \textit{life rays}.
The model begins with the idea that every person occupies their own temporal-spatial trajectory -- their own ray of life
-- and that encounters between people generate structural imprints on these rays.
These imprints persist, and they shape identity.

This essay is therefore both personal and structural. It is not only about feeling but about form.
The work should be read as a work in progress: it is intentionally open to refinement, contradiction,
and extension -- both in philosophical dialogue and in applied domains such as psychology, social dynamics, and computational systems.
